%% ECON 282E -- Homework 1, Part A
%% Reworked from PS1.lyx / PS1_IUBE.lyx (Quant Macro, HUST 2019), with the
%% recorded defects fixed: Z and beta are now stated, the Z^b typo is gone, the
%% accuracy measure is unit-free, and the step references are correct.
%% Build:  python3 tools/build_docs.py homeworks-exams/HW1/hw1a.tex
%% With answers:  python3 tools/build_docs.py --solutions homeworks-exams/HW1/hw1a.tex
%% Answers live in solutions/ (gitignored); solution PDFs are gitignored too.
%% Neither must ever be committed.

\documentclass[11pt]{article}

\newif\ifsolutions \solutionsfalse

\usepackage[margin=1in]{geometry}
\usepackage{amsmath,amssymb,booktabs,enumitem,xcolor}
\usepackage[T1]{fontenc}
\usepackage{microtype}
\usepackage[colorlinks=true,linkcolor=blue!50!black,urlcolor=blue!50!black]{hyperref}

\newcommand{\solution}[1]{\ifsolutions\medskip
  {\color{blue!45!black}\textbf{Solution.} #1}\medskip\fi}

%% The answer key lives in solutions/ (gitignored, never published), so this
%% source carries no answers. --solutions pulls it in; a plain build ignores it.
\AtBeginDocument{\ifsolutions
  \IfFileExists{solutions/hw1a_key.tex}{\input{solutions/hw1a_key.tex}}
    {\GenericWarning{}{Answer key solutions/hw1a_key.tex not found}}\fi}

\title{\vspace{-2em}\textbf{ECON 282E --- Homework 1, Part A}\\[2mm]
       \large Dynamic Programming on the Computer}
\author{Foundations of Macroeconomics \\ University of California, Riverside}
\date{Out: Thursday, October 8, 2026 \quad\textbullet\quad Due: Thursday, October 22, 2026}

\begin{document}
\maketitle
\thispagestyle{empty}

\noindent\rule{\linewidth}{0.4pt}\par
\small
\textbf{What to hand in.} One PDF with your answers and figures, the code that produced them,
and an \texttt{AI\_LOG.md} recording how you used any AI assistance. Part B is released on
October 15 and is due the same day as this part.

\textbf{On the computational questions.} Every number you report must come from code you ran.
State the grid, the tolerance and the dtype alongside each result --- the reporting standard
from Session 3. Question 2(c) uses techniques taught on \textbf{October 15}; you can do 2(a)
and 2(b) as soon as this sheet is out.

\textbf{AI policy.} Assistance is allowed and expected. You remain responsible for every
number. The validation is the part you cannot delegate.
\normalsize\par
\noindent\rule{\linewidth}{0.4pt}\par

%=============================================================================
\section*{Question 1. The model on paper (30 points)}

A representative household solves
\[
  \max_{\{c_t,k_{t+1}\}_{t=0}^{\infty}} \; \sum_{t=0}^{\infty}\beta^{t}\ln c_t
  \qquad\text{s.t.}\qquad
  c_t+k_{t+1}=Z k_t^{\alpha}, \quad c_t\ge0,\; k_{t+1}\ge0,\; k_0>0 ,
\]
with $\alpha\in(0,1)$, $\beta\in(0,1)$, $Z>0$, and \textbf{full depreciation}, $\delta=1$.
Throughout Part A use $\alpha=0.36$, $\beta=0.95$ and $Z=1$.

\begin{enumerate}[label=(\alph*), leftmargin=*]

\item \textbf{(5 pts)} Write the sequence problem and the corresponding Bellman equation.
State the state variable, the control, and the feasible correspondence $\Gamma(k)$.

\solution{\keyAone}

\item \textbf{(8 pts)} Verify Blackwell's sufficient conditions --- monotonicity and
discounting --- for the operator $T$. What do they, together with the contraction mapping
theorem, allow you to conclude about existence, uniqueness and convergence?

\solution{\keyAtwo}

\item \textbf{(10 pts)} Guess $v(k)=A+B\ln k$ and verify it. Solve for $A$ and $B$, and show
that the optimal policy is $k'=g(k)=\alpha\beta Z k^{\alpha}$ and consumption is
$c=(1-\alpha\beta)Zk^{\alpha}$. Report $A$ and $B$ numerically at the calibration above.

\solution{\keyAthree}

\item \textbf{(7 pts)} The contraction theorem gives two different bounds. Using
$\lVert v_n-v^{*}\rVert\le\frac{\beta}{1-\beta}\lVert v_n-v_{n-1}\rVert$, answer:
\begin{enumerate}[label=(\roman*)]
  \item If you stop when successive iterates differ by $\varepsilon=10^{-6}$, how far from
        $v^{*}$ might you still be, at $\beta=0.95$? At $\beta=0.99$?
  \item Using $n\gtrsim\ln\varepsilon/\ln\beta$, how many sweeps does $\varepsilon=10^{-6}$
        require at each of those two discount factors?
\end{enumerate}
Comment on what this implies for reporting a ``converged'' solution.

\solution{\keyAfour}

\end{enumerate}

%=============================================================================
\section*{Question 2. The model on the computer, parts (a)--(c) (70 points)}

Solve the same model numerically. Use a capital grid on $[0.5k^{*},1.5k^{*}]$ where
$k^{*}=(\alpha\beta Z)^{1/(1-\alpha)}$ is the deterministic steady state, and \textbf{float64}
throughout.

\begin{enumerate}[label=(\alph*), leftmargin=*]

\item \textbf{(25 pts) Grid value function iteration.} Implement VFI with the choice $k'$
restricted to the grid (no interpolation, no continuous maximization). With $N=200$ and
tolerance $10^{-6}$, report:
\begin{enumerate}[label=(\roman*)]
  \item the number of sweeps, and how it compares with your answer to 1(d)(ii);
  \item $\max_k|v(k)-v^{*}(k)|$ and the maximum \emph{relative} policy error
        $\max_k|g(k)/g^{*}(k)-1|$, against the closed form of 1(c);
  \item a plot of the computed policy against $\alpha\beta Zk^{\alpha}$, zoomed into a window
        of about 30 consecutive grid points.
\end{enumerate}
What does the zoomed plot show, and why does it not improve with more sweeps?

\solution{\keyAfive}

\item \textbf{(25 pts) Refinement, and the error budget.} Solve at $N\in\{50,100,200,400,800\}$
and tabulate the value and policy errors against $h$. Fit the slopes of $\log(\text{error})$
against $\log h$.
\begin{enumerate}[label=(\roman*)]
  \item What slopes do you get, and what do Santos \& Vigo-Aguiar (1998) predict?
  \item Repeat with tolerance $10^{-11}$. One of the two slopes changes materially. Which,
        and why? Relate your answer to 1(d)(i).
\end{enumerate}

\solution{\keyAsix}

\item \textbf{(20 pts) Acceleration.} \emph{(Techniques from Session 4, October 15.)} Starting
from your (a) solver, add each of the following and report a table of policy updates, wall
time and the maximum relative policy error:
\begin{enumerate}[label=(\roman*)]
  \item Howard's policy improvement with $n_H\in\{10,20,50\}$;
  \item exploiting monotonicity of $g$ to restrict the search;
  \item both together.
\end{enumerate}
Two things to comment on: which of these changes the \emph{answer} and which only changes the
route to it; and whether your wall-clock times move in the same direction as your operation
counts. Explain any disagreement.

\solution{\keyAseven}

\end{enumerate}

\vfill
\noindent\rule{\linewidth}{0.4pt}\par
\small
\textbf{Reminder.} Report the grid, the tolerance, the dtype and the hardware with every timing.
An accuracy claim without a benchmark, and a speed claim without a timing, both count as
unsupported.
\normalsize

\end{document}
